% =====================================================================
%  theme_tables.tex --- the one style of the project's generated tables
%
%  The highlight tint, the row-highlight command, the two-level group
%  header and the note line under a table.  Every table generated by
%  analysis/make_tables.py (the paper's Results tables) and by
%  analysis/make_report.py (the report's tables) is typeset through
%  these commands and nothing else, so the style cannot drift between
%  the two documents.
%
%  The row highlight is the light counterpart of the paper's primary
%  blue: uavblue is RGB 37,99,235, and the row tint is the same hue at
%  the light end, RGB 219,234,254, so the eye finds the row without the
%  fill fighting the type.
%
%  Needs: xcolor (for the colour), booktabs (for the rules) and colortbl
%  (for the machinery that paints a row — \rowcolor colours the
%  inter-column padding, which xcolor alone cannot).  Both documents
%  carry them in their preambles: the paper loads xcolor at
%  Mlawn-midsem-tex.tex line 10, colortbl at line 13 (added for this
%  style), booktabs at line 17; the report loads the same three and the
%  file itself, from analysis/make_report.py.  colortbl cannot be
%  loaded from a fragment, because \RequirePackage and the
%  package-opening commands are preamble-only and a generated table
%  fragment sits inside the document body — so the packages stay in the
%  preamble, and this file is written to survive a body read too: the
%  generated tables bootstrap it (the way the figure fragments
%  bootstrap analysis/theme.tex), and without colortbl the highlighted
%  rows are bold but not tinted, with the log saying so in a line that
%  starts "theme_tables:", so the wiring gap cannot be missed.
%
%  The preamble lines a document that prints these tables needs:
%
%      \usepackage{xcolor}
%      \usepackage{colortbl}
%      % =====================================================================
%  theme_tables.tex --- the one style of the project's generated tables
%
%  The highlight tint, the row-highlight command, the two-level group
%  header and the note line under a table.  Every table generated by
%  analysis/make_tables.py (the paper's Results tables) and by
%  analysis/make_report.py (the report's tables) is typeset through
%  these commands and nothing else, so the style cannot drift between
%  the two documents.
%
%  The row highlight is the light counterpart of the paper's primary
%  blue: uavblue is RGB 37,99,235, and the row tint is the same hue at
%  the light end, RGB 219,234,254, so the eye finds the row without the
%  fill fighting the type.
%
%  Needs: xcolor (for the colour), booktabs (for the rules) and colortbl
%  (for the machinery that paints a row — \rowcolor colours the
%  inter-column padding, which xcolor alone cannot).  Both documents
%  carry them in their preambles: the paper loads xcolor at
%  Mlawn-midsem-tex.tex line 10, colortbl at line 13 (added for this
%  style), booktabs at line 17; the report loads the same three and the
%  file itself, from analysis/make_report.py.  colortbl cannot be
%  loaded from a fragment, because \RequirePackage and the
%  package-opening commands are preamble-only and a generated table
%  fragment sits inside the document body — so the packages stay in the
%  preamble, and this file is written to survive a body read too: the
%  generated tables bootstrap it (the way the figure fragments
%  bootstrap analysis/theme.tex), and without colortbl the highlighted
%  rows are bold but not tinted, with the log saying so in a line that
%  starts "theme_tables:", so the wiring gap cannot be missed.
%
%  The preamble lines a document that prints these tables needs:
%
%      \usepackage{xcolor}
%      \usepackage{colortbl}
%      % =====================================================================
%  theme_tables.tex --- the one style of the project's generated tables
%
%  The highlight tint, the row-highlight command, the two-level group
%  header and the note line under a table.  Every table generated by
%  analysis/make_tables.py (the paper's Results tables) and by
%  analysis/make_report.py (the report's tables) is typeset through
%  these commands and nothing else, so the style cannot drift between
%  the two documents.
%
%  The row highlight is the light counterpart of the paper's primary
%  blue: uavblue is RGB 37,99,235, and the row tint is the same hue at
%  the light end, RGB 219,234,254, so the eye finds the row without the
%  fill fighting the type.
%
%  Needs: xcolor (for the colour), booktabs (for the rules) and colortbl
%  (for the machinery that paints a row — \rowcolor colours the
%  inter-column padding, which xcolor alone cannot).  Both documents
%  carry them in their preambles: the paper loads xcolor at
%  Mlawn-midsem-tex.tex line 10, colortbl at line 13 (added for this
%  style), booktabs at line 17; the report loads the same three and the
%  file itself, from analysis/make_report.py.  colortbl cannot be
%  loaded from a fragment, because \RequirePackage and the
%  package-opening commands are preamble-only and a generated table
%  fragment sits inside the document body — so the packages stay in the
%  preamble, and this file is written to survive a body read too: the
%  generated tables bootstrap it (the way the figure fragments
%  bootstrap analysis/theme.tex), and without colortbl the highlighted
%  rows are bold but not tinted, with the log saying so in a line that
%  starts "theme_tables:", so the wiring gap cannot be missed.
%
%  The preamble lines a document that prints these tables needs:
%
%      \usepackage{xcolor}
%      \usepackage{colortbl}
%      % =====================================================================
%  theme_tables.tex --- the one style of the project's generated tables
%
%  The highlight tint, the row-highlight command, the two-level group
%  header and the note line under a table.  Every table generated by
%  analysis/make_tables.py (the paper's Results tables) and by
%  analysis/make_report.py (the report's tables) is typeset through
%  these commands and nothing else, so the style cannot drift between
%  the two documents.
%
%  The row highlight is the light counterpart of the paper's primary
%  blue: uavblue is RGB 37,99,235, and the row tint is the same hue at
%  the light end, RGB 219,234,254, so the eye finds the row without the
%  fill fighting the type.
%
%  Needs: xcolor (for the colour), booktabs (for the rules) and colortbl
%  (for the machinery that paints a row — \rowcolor colours the
%  inter-column padding, which xcolor alone cannot).  Both documents
%  carry them in their preambles: the paper loads xcolor at
%  Mlawn-midsem-tex.tex line 10, colortbl at line 13 (added for this
%  style), booktabs at line 17; the report loads the same three and the
%  file itself, from analysis/make_report.py.  colortbl cannot be
%  loaded from a fragment, because \RequirePackage and the
%  package-opening commands are preamble-only and a generated table
%  fragment sits inside the document body — so the packages stay in the
%  preamble, and this file is written to survive a body read too: the
%  generated tables bootstrap it (the way the figure fragments
%  bootstrap analysis/theme.tex), and without colortbl the highlighted
%  rows are bold but not tinted, with the log saying so in a line that
%  starts "theme_tables:", so the wiring gap cannot be missed.
%
%  The preamble lines a document that prints these tables needs:
%
%      \usepackage{xcolor}
%      \usepackage{colortbl}
%      \input{ml_impl/analysis/theme_tables.tex}
%
%  which is what both documents now carry.
% =====================================================================

% The marker that says the theme is loaded: the generated fragments
% test it, so the file is read once, and a second read re-defines
% nothing (a second \newcommand would err).
\def\tbltheme{2026-10-21}

% ----- the colours ---------------------------------------------------
\providecolor{tblhlblue}{RGB}{219,234,254}% the highlighted row's fill

% ----- the commands the generated tables use -------------------------

% \tblHl is the highlighted row: it is expanded as the FIRST token of
% the row, and the generator pairs it with \textbf on every cell of
% that row -- the fill says "this is the row", the weight says "these
% are its values".  Which row gets it is derived from the table's own
% numbers in the generator, never entered by hand here.
\ifdefined\rowcolor
  \newcommand{\tblHl}{\rowcolor{tblhlblue}}
\else
  % Without colortbl there is no machinery that paints a row (the
  % fill has to cover the inter-column padding, and xcolor alone
  % cannot).  \tblHl is then only the weight, and the log says why:
  % the preamble line above is the thing owed.
  \typeout{** theme_tables: colortbl is not loaded, so the highlighted
   rows are bold but not tinted; see the preamble line in
   analysis/theme_tables.tex.**}
  \newcommand{\tblHl}{\relax}
\fi

% \tblGroup{span}{title} is a group header over `span' sub-columns,
% above a \cmidrule over the same span; the generator writes the
% \cmidrule with the span's own column numbers, so the two cannot go
% out of step.
\newcommand{\tblGroup}[2]{\multicolumn{#1}{c}{\textbf{#2}}}

% \tblNote{...} is the note line under a table: small and italic, one
% paragraph, and as wide as the column so it wraps instead of
% overflowing.  A note that explains a column is pointed at from that
% column's header by a ^* on the header cell.
\newcommand{\tblNote}[1]{\par\vspace{2.5pt}%
  {\parbox{\linewidth}{\scriptsize\itshape #1}}}

%
%  which is what both documents now carry.
% =====================================================================

% The marker that says the theme is loaded: the generated fragments
% test it, so the file is read once, and a second read re-defines
% nothing (a second \newcommand would err).
\def\tbltheme{2026-10-21}

% ----- the colours ---------------------------------------------------
\providecolor{tblhlblue}{RGB}{219,234,254}% the highlighted row's fill

% ----- the commands the generated tables use -------------------------

% \tblHl is the highlighted row: it is expanded as the FIRST token of
% the row, and the generator pairs it with \textbf on every cell of
% that row -- the fill says "this is the row", the weight says "these
% are its values".  Which row gets it is derived from the table's own
% numbers in the generator, never entered by hand here.
\ifdefined\rowcolor
  \newcommand{\tblHl}{\rowcolor{tblhlblue}}
\else
  % Without colortbl there is no machinery that paints a row (the
  % fill has to cover the inter-column padding, and xcolor alone
  % cannot).  \tblHl is then only the weight, and the log says why:
  % the preamble line above is the thing owed.
  \typeout{** theme_tables: colortbl is not loaded, so the highlighted
   rows are bold but not tinted; see the preamble line in
   analysis/theme_tables.tex.**}
  \newcommand{\tblHl}{\relax}
\fi

% \tblGroup{span}{title} is a group header over `span' sub-columns,
% above a \cmidrule over the same span; the generator writes the
% \cmidrule with the span's own column numbers, so the two cannot go
% out of step.
\newcommand{\tblGroup}[2]{\multicolumn{#1}{c}{\textbf{#2}}}

% \tblNote{...} is the note line under a table: small and italic, one
% paragraph, and as wide as the column so it wraps instead of
% overflowing.  A note that explains a column is pointed at from that
% column's header by a ^* on the header cell.
\newcommand{\tblNote}[1]{\par\vspace{2.5pt}%
  {\parbox{\linewidth}{\scriptsize\itshape #1}}}

%
%  which is what both documents now carry.
% =====================================================================

% The marker that says the theme is loaded: the generated fragments
% test it, so the file is read once, and a second read re-defines
% nothing (a second \newcommand would err).
\def\tbltheme{2026-10-21}

% ----- the colours ---------------------------------------------------
\providecolor{tblhlblue}{RGB}{219,234,254}% the highlighted row's fill

% ----- the commands the generated tables use -------------------------

% \tblHl is the highlighted row: it is expanded as the FIRST token of
% the row, and the generator pairs it with \textbf on every cell of
% that row -- the fill says "this is the row", the weight says "these
% are its values".  Which row gets it is derived from the table's own
% numbers in the generator, never entered by hand here.
\ifdefined\rowcolor
  \newcommand{\tblHl}{\rowcolor{tblhlblue}}
\else
  % Without colortbl there is no machinery that paints a row (the
  % fill has to cover the inter-column padding, and xcolor alone
  % cannot).  \tblHl is then only the weight, and the log says why:
  % the preamble line above is the thing owed.
  \typeout{** theme_tables: colortbl is not loaded, so the highlighted
   rows are bold but not tinted; see the preamble line in
   analysis/theme_tables.tex.**}
  \newcommand{\tblHl}{\relax}
\fi

% \tblGroup{span}{title} is a group header over `span' sub-columns,
% above a \cmidrule over the same span; the generator writes the
% \cmidrule with the span's own column numbers, so the two cannot go
% out of step.
\newcommand{\tblGroup}[2]{\multicolumn{#1}{c}{\textbf{#2}}}

% \tblNote{...} is the note line under a table: small and italic, one
% paragraph, and as wide as the column so it wraps instead of
% overflowing.  A note that explains a column is pointed at from that
% column's header by a ^* on the header cell.
\newcommand{\tblNote}[1]{\par\vspace{2.5pt}%
  {\parbox{\linewidth}{\scriptsize\itshape #1}}}

%
%  which is what both documents now carry.
% =====================================================================

% The marker that says the theme is loaded: the generated fragments
% test it, so the file is read once, and a second read re-defines
% nothing (a second \newcommand would err).
\def\tbltheme{2026-10-21}

% ----- the colours ---------------------------------------------------
\providecolor{tblhlblue}{RGB}{219,234,254}% the highlighted row's fill

% ----- the commands the generated tables use -------------------------

% \tblHl is the highlighted row: it is expanded as the FIRST token of
% the row, and the generator pairs it with \textbf on every cell of
% that row -- the fill says "this is the row", the weight says "these
% are its values".  Which row gets it is derived from the table's own
% numbers in the generator, never entered by hand here.
\ifdefined\rowcolor
  \newcommand{\tblHl}{\rowcolor{tblhlblue}}
\else
  % Without colortbl there is no machinery that paints a row (the
  % fill has to cover the inter-column padding, and xcolor alone
  % cannot).  \tblHl is then only the weight, and the log says why:
  % the preamble line above is the thing owed.
  \typeout{** theme_tables: colortbl is not loaded, so the highlighted
   rows are bold but not tinted; see the preamble line in
   analysis/theme_tables.tex.**}
  \newcommand{\tblHl}{\relax}
\fi

% \tblGroup{span}{title} is a group header over `span' sub-columns,
% above a \cmidrule over the same span; the generator writes the
% \cmidrule with the span's own column numbers, so the two cannot go
% out of step.
\newcommand{\tblGroup}[2]{\multicolumn{#1}{c}{\textbf{#2}}}

% \tblNote{...} is the note line under a table: small and italic, one
% paragraph, and as wide as the column so it wraps instead of
% overflowing.  A note that explains a column is pointed at from that
% column's header by a ^* on the header cell.
\newcommand{\tblNote}[1]{\par\vspace{2.5pt}%
  {\parbox{\linewidth}{\scriptsize\itshape #1}}}
